\newcommand{\coursemodule}{Module 12}
% ════════════════════════════════════════════════════════════════════════════
% Pixels, Frames & Transformers: shared Beamer preamble (WUR style)
% Adapted from the preamble of Bags, Vectors & Transformers.
%
% Use in a module deck (from <part>/<module>/slides/):
%   \newcommand{\coursemodule}{Module 01}
%   % ════════════════════════════════════════════════════════════════════════════
% Pixels, Frames & Transformers: shared Beamer preamble (WUR style)
% Adapted from the preamble of Bags, Vectors & Transformers.
%
% Use in a module deck (from <part>/<module>/slides/):
%   \newcommand{\coursemodule}{Module 01}
%   % ════════════════════════════════════════════════════════════════════════════
% Pixels, Frames & Transformers: shared Beamer preamble (WUR style)
% Adapted from the preamble of Bags, Vectors & Transformers.
%
% Use in a module deck (from <part>/<module>/slides/):
%   \newcommand{\coursemodule}{Module 01}
%   \input{../../../shared/preamble.tex}
%   \begin{document}
%   \moduletitle{Images as Data}
%   ...
% Compile with: latexmk -pdf <deck>.tex
% ════════════════════════════════════════════════════════════════════════════

\documentclass[aspectratio=169]{beamer}

% ── Packages ────────────────────────────────────────────────────────────────
\usepackage[T1]{fontenc}
\usepackage[utf8]{inputenc}
\usepackage{lmodern}
\usepackage{microtype}
\usepackage{xcolor}
\usepackage{tikz}
\usetikzlibrary{arrows.meta, positioning, calc, shapes.geometric, decorations.pathreplacing, matrix, fit}
\usepackage{booktabs}
\usepackage{fontawesome5}
\usepackage{array}
\usepackage[most]{tcolorbox}

\providecommand{\coursemodule}{}

% ── WUR color palette ───────────────────────────────────────────────────────
\definecolor{wurgreen}{HTML}{34B233}
\definecolor{darkgreen}{HTML}{1F6B1E}
\definecolor{lightgreen}{HTML}{E8F7E8}
\definecolor{darknavy}{HTML}{1A1A2E}
\definecolor{midgray}{HTML}{6B7280}
\definecolor{lightgray}{HTML}{F3F4F6}
\definecolor{bodytext}{HTML}{1F2937}
\definecolor{codegreen}{HTML}{9AE6B4}
\definecolor{warnred}{HTML}{FEE2E2}
\definecolor{warnborder}{HTML}{EF4444}

% ── Beamer theme ────────────────────────────────────────────────────────────
\usetheme{default}
\usecolortheme{default}
\usefonttheme{professionalfonts}
\setbeamertemplate{navigation symbols}{}
\setbeamertemplate{itemize items}{\textcolor{wurgreen}{\textbullet}}
\setbeamertemplate{itemize subitem}{\textcolor{midgray}{--}}
\setbeamertemplate{enumerate items}{\textcolor{wurgreen}{\textbf{\insertenumlabel.}}}
\setbeamercolor{frametitle}{fg=white, bg=darknavy}
\setbeamertemplate{frametitle}{%
  \begin{beamercolorbox}[wd=\paperwidth, ht=1.1cm, dp=0.2cm, leftskip=0.6cm]{frametitle}
    \usebeamerfont{frametitle}\insertframetitle
  \end{beamercolorbox}%
}
\setbeamerfont{frametitle}{size=\large, series=\bfseries}
\setbeamercolor{background canvas}{bg=white}
\setbeamercolor{normal text}{fg=bodytext}
\setbeamerfont{normal text}{size=\small}
\setbeamertemplate{footline}{%
  \begin{beamercolorbox}[wd=\paperwidth, ht=0.4cm, dp=0.15cm]{footline}%
    \hspace{0.5cm}%
    \textcolor{midgray}{\tiny Pixels, Frames \& Transformers · \coursemodule\ · Wageningen University \& Research}%
    \hfill%
    \textcolor{midgray}{\tiny \insertframenumber}%
    \hspace{0.4cm}%
  \end{beamercolorbox}%
}

% ── Boxes ───────────────────────────────────────────────────────────────────
\tcbset{wurbase/.style={arc=3pt, boxrule=0.6pt, left=7pt, right=7pt, top=5pt, bottom=5pt,
                        fontupper=\small, before skip=4pt, after skip=4pt}}
\newtcolorbox{greenbox}[1][]{wurbase, colback=lightgreen, colframe=wurgreen, #1}
\newtcolorbox{graybox}[1][]{wurbase, colback=lightgray, colframe=lightgray, #1}
\newtcolorbox{navybox}[1][]{wurbase, colback=darknavy, colframe=darknavy, colupper=white, #1}
\newtcolorbox{warnbox}[1][]{wurbase, colback=warnred, colframe=warnborder, #1}
\newtcolorbox{codebox}[1][]{wurbase, colback=darknavy, colframe=darknavy, colupper=white,
                            fontupper=\ttfamily\scriptsize, left=6pt, top=4pt, bottom=4pt, #1}

% ── Title slide ─────────────────────────────────────────────────────────────
\newcommand{\moduletitle}[1]{%
{
\setbeamertemplate{footline}{}
\begin{frame}[plain]
  \begin{tikzpicture}[remember picture, overlay]
    \fill[darknavy] (current page.south west) rectangle (current page.north east);
    \fill[wurgreen] (current page.south west) rectangle ([xshift=0.45cm]current page.north west);
    \node[anchor=west, text=white, font=\huge\bfseries, text width=13cm]
      at ([xshift=1.1cm, yshift=1.2cm]current page.west) {Pixels, Frames \& Transformers};
    \node[anchor=west, text=wurgreen, font=\large, text width=13cm]
      at ([xshift=1.1cm, yshift=0.05cm]current page.west) {#1 · \coursemodule};
    \node[anchor=west, text=midgray, font=\small, text width=12cm]
      at ([xshift=1.1cm, yshift=-1.5cm]current page.west)
      {Denise J. Roth, PhD Candidate\\
       Strategic Communication Group · Wageningen University \& Research};
  \end{tikzpicture}
\end{frame}
}}

% ── Small helpers ───────────────────────────────────────────────────────────
\newcommand{\keyterm}[1]{\textbf{\textcolor{darkgreen}{#1}}}
\newcommand{\muted}[1]{\textcolor{midgray}{#1}}
\newcommand{\credit}[1]{{\tiny\textcolor{midgray}{#1}}}
\newcommand{\notebookslide}[1]{%
\begin{frame}{In the Notebook}
\vspace{0.1cm}
#1
\end{frame}}
\newcommand{\nbitem}[2]{\item[\textcolor{wurgreen}{#1}] #2}

%   \begin{document}
%   \moduletitle{Images as Data}
%   ...
% Compile with: latexmk -pdf <deck>.tex
% ════════════════════════════════════════════════════════════════════════════

\documentclass[aspectratio=169]{beamer}

% ── Packages ────────────────────────────────────────────────────────────────
\usepackage[T1]{fontenc}
\usepackage[utf8]{inputenc}
\usepackage{lmodern}
\usepackage{microtype}
\usepackage{xcolor}
\usepackage{tikz}
\usetikzlibrary{arrows.meta, positioning, calc, shapes.geometric, decorations.pathreplacing, matrix, fit}
\usepackage{booktabs}
\usepackage{fontawesome5}
\usepackage{array}
\usepackage[most]{tcolorbox}

\providecommand{\coursemodule}{}

% ── WUR color palette ───────────────────────────────────────────────────────
\definecolor{wurgreen}{HTML}{34B233}
\definecolor{darkgreen}{HTML}{1F6B1E}
\definecolor{lightgreen}{HTML}{E8F7E8}
\definecolor{darknavy}{HTML}{1A1A2E}
\definecolor{midgray}{HTML}{6B7280}
\definecolor{lightgray}{HTML}{F3F4F6}
\definecolor{bodytext}{HTML}{1F2937}
\definecolor{codegreen}{HTML}{9AE6B4}
\definecolor{warnred}{HTML}{FEE2E2}
\definecolor{warnborder}{HTML}{EF4444}

% ── Beamer theme ────────────────────────────────────────────────────────────
\usetheme{default}
\usecolortheme{default}
\usefonttheme{professionalfonts}
\setbeamertemplate{navigation symbols}{}
\setbeamertemplate{itemize items}{\textcolor{wurgreen}{\textbullet}}
\setbeamertemplate{itemize subitem}{\textcolor{midgray}{--}}
\setbeamertemplate{enumerate items}{\textcolor{wurgreen}{\textbf{\insertenumlabel.}}}
\setbeamercolor{frametitle}{fg=white, bg=darknavy}
\setbeamertemplate{frametitle}{%
  \begin{beamercolorbox}[wd=\paperwidth, ht=1.1cm, dp=0.2cm, leftskip=0.6cm]{frametitle}
    \usebeamerfont{frametitle}\insertframetitle
  \end{beamercolorbox}%
}
\setbeamerfont{frametitle}{size=\large, series=\bfseries}
\setbeamercolor{background canvas}{bg=white}
\setbeamercolor{normal text}{fg=bodytext}
\setbeamerfont{normal text}{size=\small}
\setbeamertemplate{footline}{%
  \begin{beamercolorbox}[wd=\paperwidth, ht=0.4cm, dp=0.15cm]{footline}%
    \hspace{0.5cm}%
    \textcolor{midgray}{\tiny Pixels, Frames \& Transformers · \coursemodule\ · Wageningen University \& Research}%
    \hfill%
    \textcolor{midgray}{\tiny \insertframenumber}%
    \hspace{0.4cm}%
  \end{beamercolorbox}%
}

% ── Boxes ───────────────────────────────────────────────────────────────────
\tcbset{wurbase/.style={arc=3pt, boxrule=0.6pt, left=7pt, right=7pt, top=5pt, bottom=5pt,
                        fontupper=\small, before skip=4pt, after skip=4pt}}
\newtcolorbox{greenbox}[1][]{wurbase, colback=lightgreen, colframe=wurgreen, #1}
\newtcolorbox{graybox}[1][]{wurbase, colback=lightgray, colframe=lightgray, #1}
\newtcolorbox{navybox}[1][]{wurbase, colback=darknavy, colframe=darknavy, colupper=white, #1}
\newtcolorbox{warnbox}[1][]{wurbase, colback=warnred, colframe=warnborder, #1}
\newtcolorbox{codebox}[1][]{wurbase, colback=darknavy, colframe=darknavy, colupper=white,
                            fontupper=\ttfamily\scriptsize, left=6pt, top=4pt, bottom=4pt, #1}

% ── Title slide ─────────────────────────────────────────────────────────────
\newcommand{\moduletitle}[1]{%
{
\setbeamertemplate{footline}{}
\begin{frame}[plain]
  \begin{tikzpicture}[remember picture, overlay]
    \fill[darknavy] (current page.south west) rectangle (current page.north east);
    \fill[wurgreen] (current page.south west) rectangle ([xshift=0.45cm]current page.north west);
    \node[anchor=west, text=white, font=\huge\bfseries, text width=13cm]
      at ([xshift=1.1cm, yshift=1.2cm]current page.west) {Pixels, Frames \& Transformers};
    \node[anchor=west, text=wurgreen, font=\large, text width=13cm]
      at ([xshift=1.1cm, yshift=0.05cm]current page.west) {#1 · \coursemodule};
    \node[anchor=west, text=midgray, font=\small, text width=12cm]
      at ([xshift=1.1cm, yshift=-1.5cm]current page.west)
      {Denise J. Roth, PhD Candidate\\
       Strategic Communication Group · Wageningen University \& Research};
  \end{tikzpicture}
\end{frame}
}}

% ── Small helpers ───────────────────────────────────────────────────────────
\newcommand{\keyterm}[1]{\textbf{\textcolor{darkgreen}{#1}}}
\newcommand{\muted}[1]{\textcolor{midgray}{#1}}
\newcommand{\credit}[1]{{\tiny\textcolor{midgray}{#1}}}
\newcommand{\notebookslide}[1]{%
\begin{frame}{In the Notebook}
\vspace{0.1cm}
#1
\end{frame}}
\newcommand{\nbitem}[2]{\item[\textcolor{wurgreen}{#1}] #2}

%   \begin{document}
%   \moduletitle{Images as Data}
%   ...
% Compile with: latexmk -pdf <deck>.tex
% ════════════════════════════════════════════════════════════════════════════

\documentclass[aspectratio=169]{beamer}

% ── Packages ────────────────────────────────────────────────────────────────
\usepackage[T1]{fontenc}
\usepackage[utf8]{inputenc}
\usepackage{lmodern}
\usepackage{microtype}
\usepackage{xcolor}
\usepackage{tikz}
\usetikzlibrary{arrows.meta, positioning, calc, shapes.geometric, decorations.pathreplacing, matrix, fit}
\usepackage{booktabs}
\usepackage{fontawesome5}
\usepackage{array}
\usepackage[most]{tcolorbox}

\providecommand{\coursemodule}{}

% ── WUR color palette ───────────────────────────────────────────────────────
\definecolor{wurgreen}{HTML}{34B233}
\definecolor{darkgreen}{HTML}{1F6B1E}
\definecolor{lightgreen}{HTML}{E8F7E8}
\definecolor{darknavy}{HTML}{1A1A2E}
\definecolor{midgray}{HTML}{6B7280}
\definecolor{lightgray}{HTML}{F3F4F6}
\definecolor{bodytext}{HTML}{1F2937}
\definecolor{codegreen}{HTML}{9AE6B4}
\definecolor{warnred}{HTML}{FEE2E2}
\definecolor{warnborder}{HTML}{EF4444}

% ── Beamer theme ────────────────────────────────────────────────────────────
\usetheme{default}
\usecolortheme{default}
\usefonttheme{professionalfonts}
\setbeamertemplate{navigation symbols}{}
\setbeamertemplate{itemize items}{\textcolor{wurgreen}{\textbullet}}
\setbeamertemplate{itemize subitem}{\textcolor{midgray}{--}}
\setbeamertemplate{enumerate items}{\textcolor{wurgreen}{\textbf{\insertenumlabel.}}}
\setbeamercolor{frametitle}{fg=white, bg=darknavy}
\setbeamertemplate{frametitle}{%
  \begin{beamercolorbox}[wd=\paperwidth, ht=1.1cm, dp=0.2cm, leftskip=0.6cm]{frametitle}
    \usebeamerfont{frametitle}\insertframetitle
  \end{beamercolorbox}%
}
\setbeamerfont{frametitle}{size=\large, series=\bfseries}
\setbeamercolor{background canvas}{bg=white}
\setbeamercolor{normal text}{fg=bodytext}
\setbeamerfont{normal text}{size=\small}
\setbeamertemplate{footline}{%
  \begin{beamercolorbox}[wd=\paperwidth, ht=0.4cm, dp=0.15cm]{footline}%
    \hspace{0.5cm}%
    \textcolor{midgray}{\tiny Pixels, Frames \& Transformers · \coursemodule\ · Wageningen University \& Research}%
    \hfill%
    \textcolor{midgray}{\tiny \insertframenumber}%
    \hspace{0.4cm}%
  \end{beamercolorbox}%
}

% ── Boxes ───────────────────────────────────────────────────────────────────
\tcbset{wurbase/.style={arc=3pt, boxrule=0.6pt, left=7pt, right=7pt, top=5pt, bottom=5pt,
                        fontupper=\small, before skip=4pt, after skip=4pt}}
\newtcolorbox{greenbox}[1][]{wurbase, colback=lightgreen, colframe=wurgreen, #1}
\newtcolorbox{graybox}[1][]{wurbase, colback=lightgray, colframe=lightgray, #1}
\newtcolorbox{navybox}[1][]{wurbase, colback=darknavy, colframe=darknavy, colupper=white, #1}
\newtcolorbox{warnbox}[1][]{wurbase, colback=warnred, colframe=warnborder, #1}
\newtcolorbox{codebox}[1][]{wurbase, colback=darknavy, colframe=darknavy, colupper=white,
                            fontupper=\ttfamily\scriptsize, left=6pt, top=4pt, bottom=4pt, #1}

% ── Title slide ─────────────────────────────────────────────────────────────
\newcommand{\moduletitle}[1]{%
{
\setbeamertemplate{footline}{}
\begin{frame}[plain]
  \begin{tikzpicture}[remember picture, overlay]
    \fill[darknavy] (current page.south west) rectangle (current page.north east);
    \fill[wurgreen] (current page.south west) rectangle ([xshift=0.45cm]current page.north west);
    \node[anchor=west, text=white, font=\huge\bfseries, text width=13cm]
      at ([xshift=1.1cm, yshift=1.2cm]current page.west) {Pixels, Frames \& Transformers};
    \node[anchor=west, text=wurgreen, font=\large, text width=13cm]
      at ([xshift=1.1cm, yshift=0.05cm]current page.west) {#1 · \coursemodule};
    \node[anchor=west, text=midgray, font=\small, text width=12cm]
      at ([xshift=1.1cm, yshift=-1.5cm]current page.west)
      {Denise J. Roth, PhD Candidate\\
       Strategic Communication Group · Wageningen University \& Research};
  \end{tikzpicture}
\end{frame}
}}

% ── Small helpers ───────────────────────────────────────────────────────────
\newcommand{\keyterm}[1]{\textbf{\textcolor{darkgreen}{#1}}}
\newcommand{\muted}[1]{\textcolor{midgray}{#1}}
\newcommand{\credit}[1]{{\tiny\textcolor{midgray}{#1}}}
\newcommand{\notebookslide}[1]{%
\begin{frame}{In the Notebook}
\vspace{0.1cm}
#1
\end{frame}}
\newcommand{\nbitem}[2]{\item[\textcolor{wurgreen}{#1}] #2}


\begin{document}

\moduletitle{Detection, Segmentation and Counting}

% ─────────────────────────────────────────────────────────────────────────────
\begin{frame}{What Is in the Image, and Where?}
\begin{columns}[T]
\begin{column}{0.31\textwidth}
\begin{graybox}[height=3.2cm]
\raggedright\textbf{Detection}\\
boxes around objects, each with a category and a confidence score
\end{graybox}
\end{column}
\begin{column}{0.31\textwidth}
\begin{greenbox}[height=3.2cm]
\raggedright\textbf{Segmentation}\\
the exact outline of an object, pixel by pixel
\end{greenbox}
\end{column}
\begin{column}{0.31\textwidth}
\begin{navybox}[height=3.2cm]
\raggedright\textbf{Counting}\\
how many of something are visible
\end{navybox}
\end{column}
\end{columns}

\vspace{0.25cm}
They let you code \emph{parts} of an image: how many people are at a rally, whether police appear,
how much of a cereal box a mascot takes up.
\end{frame}

% ─────────────────────────────────────────────────────────────────────────────
\begin{frame}{Standard Detectors Know 80 Things}
\begin{columns}[T]
\begin{column}{0.5\textwidth}
\textbf{RT-DETR}, a transformer-based detector, is trained on \textbf{COCO}: 80 everyday
categories such as \emph{person}, \emph{car}, \emph{umbrella}, \emph{backpack}.

\vspace{0.15cm}
No \emph{police officer}, no \emph{banner}, no \emph{flag}. Like ImageNet's categories (module 06),
they were chosen for a benchmark, not for your question.
\end{column}
\begin{column}{0.46\textwidth}
\begin{graybox}
\textbf{The confidence threshold}\\
Every box has a score. A low threshold finds more objects and more false ones; a high threshold
misses partly hidden ones. Keep all boxes above a low score once, then try thresholds without
rerunning the model.
\end{graybox}
\end{column}
\end{columns}
\end{frame}

% ─────────────────────────────────────────────────────────────────────────────
\begin{frame}{Counting People Is Hard}
\begin{columns}[T]
\begin{column}{0.5\textwidth}
\begin{itemize}
\item People hide behind each other
\item Distant heads are a few pixels tall
\item Detectors return a limited number of boxes
\item The frame shows one moment and one part of the event
\end{itemize}

\vspace{0.15cm}
\textbf{Validation}: your codebook says \emph{crowd} = ten or more visible people, so ``10+
detected'' is a rule fixed in advance that can be scored on the test set.
\end{column}
\begin{column}{0.46\textwidth}
\begin{warnbox}
\textbf{A count is a lower bound} on the people visible and detectable in one photo, not the size
of a demonstration.
\end{warnbox}
\muted{\footnotesize Crowd size estimation uses area times density, often from several or aerial
images, and reports a range.}
\end{column}
\end{columns}
\end{frame}

% ─────────────────────────────────────────────────────────────────────────────
\begin{frame}{Open-Vocabulary Detection}
\begin{columns}[T]
\begin{column}{0.5\textwidth}
\textbf{OWLv2} scores every candidate box against \textbf{text queries} with a CLIP-style text
encoder (module 09):

\vspace{0.1cm}
{\footnotesize\texttt{"a police officer", "a protest sign", "a banner", "a flag"}}

\vspace{0.15cm}
Image-level codes: the highest score per query, above a threshold, validated against your
\emph{police} and \emph{signs} codes.
\end{column}
\begin{column}{0.46\textwidth}
\begin{warnbox}
\textbf{Threshold and wording both matter.} ``police'', ``a police officer in uniform'' and ``a
person wearing a police uniform'' give different results.
\end{warnbox}
\begin{greenbox}
Choose them on a separate hand-coded validation sample, then score the test set once.
\end{greenbox}
\end{column}
\end{columns}
\end{frame}

% ─────────────────────────────────────────────────────────────────────────────
\begin{frame}{Segmentation with SAM 2}
\begin{columns}[T]
\begin{column}{0.5\textwidth}
\textbf{SAM 2} \muted{(Meta, Apache 2.0)} outlines an object when you point at it with a click or a
box. It doesn't know \emph{what} it outlines, so we combine it:

\vspace{0.1cm}
\begin{enumerate}
\item OWLv2 finds ``a cartoon character''
\item SAM 2 traces it inside the box
\item measure: the share of the package front covered
\end{enumerate}
\end{column}
\begin{column}{0.46\textwidth}
\centering
\begin{tikzpicture}
\fill[lightgray] (0, 0) rectangle (3.6, 3.0);
\fill[wurgreen!70] (1.8, 1.5) +(90:1.1) -- +(162:0.45) -- +(162:1.1) -- +(234:0.45) -- +(234:1.1)
  -- +(306:0.45) -- +(306:1.1) -- +(18:0.45) -- +(18:1.1) -- +(90:0.45) -- cycle;
\draw[yellow!80!black, very thick] (0.75, 0.6) rectangle (2.85, 2.6);
\node[font=\scriptsize, align=center] at (1.8, -0.35) {here the box is almost three times the area of the mask};
\end{tikzpicture}
\end{column}
\end{columns}

\vspace{0.15cm}
\muted{\footnotesize Kirillov et al.\ (2023) introduced Segment Anything; SAM 2 extends it to video
(Ravi et al., 2024).}
\end{frame}

% ─────────────────────────────────────────────────────────────────────────────
\begin{frame}{Measures Built from Regions}
\begin{columns}[T]
\begin{column}{0.5\textwidth}
\begin{itemize}
\item \keyterm{count}: people, flags, police vehicles
\item \keyterm{size}: share of the image a mascot or a speaker takes up
\item \keyterm{position}: foreground or background, center or edge
\item \keyterm{overlap}: is the sign held by a person?
\end{itemize}
\end{column}
\begin{column}{0.46\textwidth}
\begin{graybox}
\textbf{Before you claim} that mascots are used to sell less healthy cereals, validate the detector
against hand codes, check that a few brands with many variants don't drive the pattern, and think
about confounders such as the target age group.
\end{graybox}
\end{column}
\end{columns}
\end{frame}

% ─────────────────────────────────────────────────────────────────────────────
\begin{frame}{Two-Level Classification}
\vspace{0.1cm}
\centering
\begin{tikzpicture}[
  b/.style={draw=midgray, fill=lightgray, rounded corners=3pt, font=\footnotesize, align=center, inner sep=5pt},
  g/.style={draw=wurgreen, fill=lightgreen, rounded corners=3pt, font=\footnotesize, align=center, inner sep=5pt},
  n/.style={draw=darknavy, fill=darknavy, text=white, rounded corners=3pt, font=\footnotesize, align=center, inner sep=5pt},
  arr/.style={-{Stealth[length=2mm]}, thick, color=midgray}]
\node[b, text width=1.6cm] (img) at (0, 0) {image};
\node[g, text width=3.2cm] (parts) at (3.6, 0) {\textbf{level 1}: detected parts\\ people, cars, signs, police, flags};
\node[g, text width=2.8cm] (model) at (7.9, 0) {\textbf{level 2}: a logistic regression on the parts};
\node[n, text width=1.6cm] (out) at (11.4, 0) {protest or not};
\draw[arr] (img) -- (parts);
\draw[arr] (parts) -- (model);
\draw[arr] (model) -- (out);
\end{tikzpicture}

\vspace{0.4cm}
\raggedright
\textbf{Readable}: the coefficients say which objects push an image toward ``protest'', and expose
shortcuts (a large weight on \emph{umbrella} says something about the weather in your training
data). \textbf{The price}: whatever the parts don't capture (setting, mood) is lost. Adding CLIP
embeddings usually scores higher and reads worse.

\vspace{0.15cm}
\muted{\footnotesize Scholz, Weidmann, Steinert-Threlkeld, Keremoglu \& Goldlücke (2025),
\emph{Political Analysis}.}
\end{frame}

% ─────────────────────────────────────────────────────────────────────────────
\begin{frame}{Licenses: Check the Code and the Weights}
\begin{columns}[T]
\begin{column}{0.5\textwidth}
Many detection tutorials use YOLO models from the Ultralytics library. Its recent versions are
licensed under \textbf{AGPL-3.0}: software built with it must be released under the same license,
also when offered as a web service, unless you buy a commercial license.
\end{column}
\begin{column}{0.46\textwidth}
\begin{greenbox}
\textbf{In this module}: RT-DETR, OWLv2 and SAM 2 are Apache 2.0, which asks only for attribution
and a copy of the license.
\end{greenbox}
\muted{\footnotesize Fine for a published research script either way, but it matters for tools
you build for partners. Code and weights can have different licenses.}
\end{column}
\end{columns}
\end{frame}

% ─────────────────────────────────────────────────────────────────────────────
\notebookslide{
\begin{itemize}
\setlength{\itemsep}{5pt}
\nbitem{\faUsers}{Detecting and counting people with RT-DETR, checked against your crowd codes}
\nbitem{\faSearch}{Finding police and signs with OWLv2, and how thresholds and wording matter}
\nbitem{\faDrawPolygon}{Outlining cereal mascots with SAM 2 and measuring how much of the box they cover}
\nbitem{\faProjectDiagram}{A two-level classifier from detected objects, and what it costs in accuracy}
\end{itemize}

\vspace{0.35cm}
\muted{\footnotesize part3\_transformers/12\_detection\_segmentation\_and\_counting/notebooks/\\
12\_detection\_segmentation\_and\_counting\_exercises.ipynb · 5 exercises · GPU needed}
}

\end{document}
